\chapter*{Acknowledgments}
\addcontentsline{toc}{chapter}{Acknowledgments}

The development of the Dialectical Cognition Framework and this treatise benefited from multiple sources of insight and support.

\textbf{Intellectual Foundations}: The philosophical traditions of Socratic inquiry, the cognitive science of extended mind, and the emerging field of human-AI interaction provided essential grounding. Thinkers from Plato to Andy Clark to the researchers at Anthropic have shaped the conceptual landscape in which DCF operates.

\textbf{Tool Creators}: The engineers and designers who built Claude, Claude Code, and the broader ecosystem of AI development tools made this work possible. Their decisions about features like Plan Mode, memory systems, and approval workflows directly influenced DCF's evolution.

\textbf{The Practice Community}: Practitioners exploring human-AI collaboration in forums, workplaces, and open source projects contributed countless insights through their experiments and shared experiences.

\textbf{The Collaboration Itself}: This document emerged through the very process it describes. The recursive refinement across multiple sessions, the Socratic questioning of assumptions, and the iterative synthesis of ideas demonstrated DCF's validity through practice rather than mere assertion.

\textbf{Future Contributors}: DCF is designed to evolve. Those who will adapt, critique, and extend these ideas are acknowledged in advance for the improvements they will make.

% ============================================================================
% COLOPHON
% ============================================================================

\chapter*{Colophon}
\addcontentsline{toc}{chapter}{Colophon}

This treatise was developed through the very methodology it describes: iterative dialogue with AI, recursive refinement of ideas, and the synthesis of multiple sources into coherent understanding. It is itself an artifact of DCF in practice.

The Dialectical Cognition Framework continues to evolve. Document your adaptations. Share what works. Build what you envision.

\begin{center}
\textbf{The architecture of thought awaits its architects.}
\end{center}
